\paragraph{補助三次曲線の基底不変性.}
$Q$ の係数の偏極から得られる $G$ は、二つのベクトル空間の基底変更に対して自然である。
$K_6$ は、$P$ の係数空間内に値を持つ一次の直交syzygy全体として定義でき、
三つの交代形式は、その二つのsyzygyの係数のpairingを平面座標について反対称化したものである。
従って $\operatorname{GL}_2$ による合同変換と平面座標の変更は、
$K_6$ の基底変更および三つの形式の線形な基底変更を誘導する。
Pfaffianの零点集合は、前者で非零スカラー倍、後者で射影座標変更を受けるだけである。
$E_Q$ および $\mathbb P_{E_Q}(\ker\Omega)$ の同型類はこれらの選択に依存しない。
三つの形式のパラメータ平面は $\mathbb P(\bigwedge^2V)$ であり、
$\dim V=3$ から自然に元の平面の双対平面 $\mathbb P(V^*)$ と同一視される。

特に、一般の平面四次Prym対 $(C,\eta)$ に対する二次対称行列表示の合同変換を除く一意性により、
$E_Q$ の同型類は $(C,\eta)$ の不変量となる。
これは、指定反対合を忘れた元の二十一点配置から、
その指定を使わずに同じ不変量を復元したという主張ではない。

\paragraph{補助楕円曲線の非定数変動.}
上記のHesse型の族に対し、通常の $j$ 不変量 $j(i)=1728$ の正規化で
\[
 j(E_t)=-\frac{(t-8)^3\bigl((t-8)^3-432t^2\bigr)^3}
 {2t^2\bigl((t-8)^3+54t^2\bigr)^3}
 =-\frac{(t-8)^3\bigl((t-8)^3-432t^2\bigr)^3}
 {2t^2(t-2)^3(t+16)^6}.
\]
これはHesse形 $X^3+Y^3+Z^3+3hXYZ=0$ の公式
$27h^3(8-h^3)^3/(1+h^3)^3$ に
$h^3=(t-8)^3/(54t^2)$ を代入して得る
\cite[\S2]{ArtebaniDolgachevHesse}。
右辺は $t=\infty$ に一位の極を持ち、従って非定数である。
よって補助曲線を与える有理写像は種数一曲線の粗モジュライ空間へ支配的である。
この結論は十八次元の行列空間上の有理関数の非定数性であるから、
元の四次Gram曲線が滑らかな稠密開部分へ制限しても変わらない。
同型不変性により、指定Prym対のモジュライ上にも非定数の不変量として降りる。
